\section*{Straight-line programs and circuit values (Appendix B)}

Take a word $w=a_0\cdots a_{n-1}\in\Sigma^*$ over the alphabet $\Sigma$ ($n\ge0$, $a_0,\ldots,a_{n-1}\in\Sigma$). For $0\le i<n$ let $w[i]=a_i$ and for $0\le i\le j<n$ let $w[i:j]=a_ia_{i+1}\cdots a_j$. Moreover $w[:i]=w[0:i]$. Note that in the notations $w[i]$ and $w[i:j]$ we take $0$ as the first position in $w$. Finally, with $|w|_a=|\{i\mid a_i=a\}|$ we denote the number of occurrences of $a$ in $w$.
A {\em straight-line program} (SLP for short) over the alphabet $\Sigma$ is a triple $\mathcal{G} = (V,\rho,S)$, where 
$V$ is a finite set of variables such that $V \cap \Sigma = \emptyset$, $S \in V$ is the start variable, and $\rho: V \to (V \cup \Sigma)^*$ is a mapping
such that the relation $\{ (A,B) \in V \times V \colon B \text{ occurs in } \rho(A) \}$
is acyclic. For the reader familiar with context free grammars, it might be helpful to view the SLP
$\mathcal{G} = (V,\rho,S)$ as the context-free grammar $(V,\Sigma, P, S)$, where $P$ contains all productions $A \to \rho(A)$ for $A \in V$.
The definition of an SLP implies that this context-free grammar derives exactly on terminal word, which will be denoted by
$\val(\mathcal{G})$. Formally, one can extend $\rho$ to a morphism $\rho : (V \cup \Sigma)^* \to (V \cup \Sigma)^*$ by setting
$\rho(a)=a$ for all $a \in \Sigma$. The above acyclicity condition on $\rho$ implies that for $m = |V|$ we have $\rho^m(w) \in \Sigma^*$ for all
$w \in (V \cup \Sigma)^*$. We then define $\val_{\mathcal{G}}(w) = \rho^m(w)$ (the string derived from the sentential form $w$) and 
$\val(\mathcal{G}) = \val_{\mathcal{G}}(S)$.

The word $\rho(A)$ is also called the {\em right-hand side} of $A$. 
Quite often, it is convenient to assume that all right-hand sides are of the form $a \in \Sigma$ or $BC$ with $B,C \in V$.
This corresponds to the well-known Chomsky normal form for context-free grammars. There is a simple linear time 
algorithm that transforms an SLP $\mathcal{G}$ with $\val(\mathcal{G}) \neq \varepsilon$ into an SLP $\mathcal{G}'$ in Chomsky normal
form with $\val(\mathcal{G}) = \val(\mathcal{G}')$, see e.g. \cite[Proposition 3.8]{ref38}.

We define the size of the SLP $\mathcal{G} = (V,\rho,S)$ as the total length of all right-hand sides:  $|\mathcal{G}| = \sum_{A \in V} |\rho(A)|$.
SLPs offer a succinct representation of words that contain many repeated substrings. For instance, the word $(ab)^{2^n}$ can be produced
by the SLP $\mathcal{G} = (\{A_0, \ldots,A_n\},\rho,A_n)$ with $\rho(A_0) = ab$ and $\rho(A_{i+1}) = A_i A_i$ for $0 \leq i \leq n-1$.
We need the following upper bound on the length of the word $\val(\mathcal{G})$:

\begin{authorclaim}{Lemma 31 (c.f.\ [13])} \label{lemma-SLP-upper-bound}
For every SLP $\mathcal{G}$ we have $|\val(\mathcal{G})| \leq  3^{|\mathcal{G}|/3}$.
\end{authorclaim}
We also need polynomial time algorithms for a few algorithmic problems for SLPs:

\begin{authorclaim}{Lemma 32 ([38, Chapter 3])} \label{lemma-basic}
The following problems can be solved in polynomial time, where $\mathcal{G}$ is an SLP over a terminal alphabet $\Sigma$, $a \in \Sigma$, and $p, q \in \mathbb{N}$
(the latter are given in binary notation):
\begin{itemize}
\item Given $\mathcal{G}$, compute the length $|\val(\mathcal{G})|$.
\item Given $\mathcal{G}$ and $a$, compute the number $|\val(\mathcal{G})|_a$ of occurrences of $a$.
\item Given $\mathcal{G}$ and $p$, compute the symbol $\val(\mathcal{G})[p] \in \Sigma$ (in case $0 \le p < |\val(\mathcal{G})|$ does not hold, the algorithm outputs a special symbol).
\item Given $\mathcal{G}$ and $p,q$, compute an SLP for the string $\val(\mathcal{G})[p:q]$ (in case $0 \le p \leq q < |\val(\mathcal{G})|$ does not hold, the algorithm
outputs a special symbol).
\end{itemize}
\end{authorclaim}

\begin{authorclaim}{Lemma 33 (c.f.\ [38, Lemma 3.12])}  \label{lemma-SLP-iterated-morph}
Given a symbols $a_0 \in \Sigma$ and a sequence of morphisms $\phi_1, \ldots, \phi_n : \Sigma^* \to \Sigma^*$, where every $\phi_i$ is given by a list of the words
$\phi_i(a)$ for $a \in \Sigma$, one can compute in \L an SLP for the word
$\phi_1(\phi_2( \cdots \phi_n(a_0) \cdots))$.
\end{authorclaim}
The {\em circuit value problem} for a finitely generated group $G$ with the finite standard generating set $\Sigma$, $\CVP(G,\Sigma)$ for short,
is the following decision problem:
\begin{description}
\item[Input:]  an SLP $\mathcal{G}$ over the alphabet $\Sigma$.
\item [Question:] does $\val(\mathcal{G})=1$ hold in $G$?
\end{description}
It is an easy observation that the computational complexity of the compressed word 
problem for $G$ does not depend on the chosen generating set $\Sigma$ in the sense
that if $\Sigma'$ is another finite standard generating set for $G$, then $\CVP(G,\Sigma)$
is \L-reducible to $\CVP(G,\Sigma')$ \cite[Lemma~4.2]{ref38}. Therefore we do not
have to specify the generating set and we just write $\CVP(G)$.

The circuit value problem for $G$ is equivalent to the problem whether a given circuit over the group $G$ evaluates to $1$: Take an SLP $\mathcal{G} = (V,\rho,S)$ in
Chomsky normal form and built a circuit by taking $V$ is the set of gates. If $\rho(A)=a  \in \Sigma$ then $A$ is an input gate that is labelled
with the group generator $a$. If $\rho(A) = BC$ with $B,C \in V$ then $B$ is left input gate for $A$ and $C$ is the right input gate for $A$.
Such a circuit can be evaluated in the natural way (every internal gate computes the product of its input values) and the circuit output is the value
at gate $S$. 


